\subsection{AC 自动机 (Aho-Corasick)}
字符集过大时，暴力跳 fail 是对的。
\begin{lstlisting}
ll ch[1000010][26];
ll fail[1000010];
vi e[1000010];
void add(ll nd){ll u=1;
	ll len=strlen(c+1);
	f(i,1,len){
		ll p=c[i]-'a';
		if(!ch[u][p])ch[u][p]=++cnt;
		u=ch[u][p];
	}
}queue<ll>q;
void gtf(){
	fail[1]=1;
	f(i,0,25){
		if(ch[1][i])q.push(ch[1][i]),fail[ch[1][i]]=1;
		else ch[1][i]=1;
	}while(!q.empty()){
		ll u=q.front();q.pop();
		f(i,0,25){
			if(ch[u][i])fail[ch[u][i]]=ch[fail[u]][i],q.push(ch[u][i]);
			else ch[u][i]=ch[fail[u]][i];
		}
	}
}

//字符集过大时，暴力跳 fail 是对的。用 find(u,i) 代替 ch[u][i]（getfail里和跑acam的时候）
ll find(ll u,ll i){
	while(u!=1&&!ch[u][i])u=fail[u];
    return ch[u][i]?ch[u][i]:1;
}
void init(){
    f(i,0,cnt){fill(ch[i],ch[i]+26,0);fail[i]=0;e[i].clear();}
    cnt=1;while(!q.empty())q.pop();
}
\end{lstlisting}

\paragraph{多测与注意} 每组先 init，再全部 add、gtf；保留根编号 1。gtf 补全转移后不要再插模式或重复构建。稀疏版 find 只用于未补全的 trie 转移。
